%%%%%%%%%%%%%%%%%%%%%%%%%%%%%%%%%%%%%%%%%
% University Assignment Title Page 
% LaTeX Template
% Version 1.0 (27/12/12)
%
% This template has been downloaded from:
% http://www.LaTeXTemplates.com
%
% Original author:
% WikiBooks (http://en.wikibooks.org/wiki/LaTeX/Title_Creation)
%
% License:
% CC BY-NC-SA 3.0 (http://creativecommons.org/licenses/by-nc-sa/3.0/)
% 
% Instructions for using this template:
% This title page is capable of being compiled as is. This is not useful for 
% including it in another document. To do this, you have two options: 
%
% 1) Copy/paste everything between \begin{document} and \end{document} 
% starting at \begin{titlepage} and paste this into another LaTeX file where you 
% want your title page.
% OR
% 2) Remove everything outside the \begin{titlepage} and \end{titlepage} and 
% move this file to the same directory as the LaTeX file you wish to add it to. 
% Then add \input{./title_page_1.tex} to your LaTeX file where you want your
% title page.
%
%%%%%%%%%%%%%%%%%%%%%%%%%%%%%%%%%%%%%%%%%
%\title{Title page with logo}
%----------------------------------------------------------------------------------------
%	PACKAGES AND OTHER DOCUMENT CONFIGURATIONS
%----------------------------------------------------------------------------------------

\documentclass[12pt]{article}
\usepackage[english]{babel}
\usepackage[utf8x]{inputenc}
\usepackage{amsmath}
\usepackage{xkeyval}
\usepackage{graphicx}
\usepackage[colorinlistoftodos]{todonotes}

\begin{document}

\begin{titlepage}

\newcommand{\HRule}{\rule{\linewidth}{0.5mm}} % Defines a new command for the horizontal lines, change thickness here

\center % Center everything on the page
 
%----------------------------------------------------------------------------------------
%	HEADING SECTIONS
%----------------------------------------------------------------------------------------

\textsc{\LARGE Pontifícia Universidade Católica Campinas}\\[1.5cm] % Name of your university/college
\textsc{\Large Engenharia de Computação}\\[0.5cm] % Major heading such as course name
\textsc{\large Sistemas Operácionais B}\\[0.5cm] % Minor heading such as course title

%----------------------------------------------------------------------------------------
%	TITLE SECTION
%----------------------------------------------------------------------------------------

\HRule \\[0.4cm]
{ \huge \bfseries Relatório Experimento 1}\\[0.4cm] % Title of your document
\HRule \\[1.5cm]
 
%----------------------------------------------------------------------------------------
%	AUTHOR SECTION
%----------------------------------------------------------------------------------------

\begin{minipage}{0.4\textwidth}
\begin{flushleft} \large
\emph{Author:}\\
Bruno Camilo \textsc{Silvério}
Guilherme \textsc{Pernicone}
João Pedro \textsc{Porta}
Marcelo Dib \textsc{Coutinho}
Pedro Pierina \textsc{Pierina}
\end{flushleft}
\end{minipage}
~
\begin{minipage}{0.4\textwidth}
\begin{flushright} \large
\emph{Supervisor:} \\
Ms. Prof. Edmar Roberto Santana de \textsc{Resende} % Supervisor's Name
\end{flushright}
\end{minipage}\\[2cm]

% If you don't want a supervisor, uncomment the two lines below and remove the section above
%\Large \emph{Author:}\\
%John \textsc{Smith}\\[3cm] % Your name

%----------------------------------------------------------------------------------------
%	DATE SECTION
%----------------------------------------------------------------------------------------

{\large \today}\\[2cm] % Date, change the \today to a set date if you want to be precise

%----------------------------------------------------------------------------------------
%	LOGO SECTION
%----------------------------------------------------------------------------------------

\includegraphics[width=26mm]{logo.jpg}% Include a department/university logo - this will require the graphicx package
 
%----------------------------------------------------------------------------------------

\vfill % Fill the rest of the page with whitespace

\end{titlepage}


\begin{abstract}
Experimento com o objetivo de aprender e entender as dificuldades de criar e colocar em uso um modulo de Kernel do sitema operácional Ubuntu 16.04.
\end{abstract}

\section{Introdução}

Nesse experimento, utilizando Ubuntu 16.04, foi criado um módulo que tem como objetivo utilizar a biblioteca crypto.h para cryptografar, decriptografar uma string dada pelo usuário atravéz de um programa de usuário, escrito em C, o programa de usuário comunica com o módulo atravéz da escrita e leitura no modulo /dev/cryptomodule 

\section{Fundamentação Teórica}

\subsection{Sistema Operacional (SO)}

Um sistema operacional (SO) funciona como uma interface entre os usuários e o hardware, ele é um programa cuja a função é gerenciar os recursos do sistema, definindo quais programas devem receber a atenção do processador e é responsável por criar um sistema de arquivos a fim de indexar a memória. Além disso o SO também é responsável por fazer a comunicação entre os programas, periféricos e o hardware da máquina. 

\subsection{Kernel}

O kernel ou núcleo é o componente central do sistema operacional. Ele serve como ponte entre os aplicativos e o processamento real dos dados feito a nível de hardware. O kernel é responsável por gerenciar os recursos dos sistemas e oferecendo ajuda para que os aplicativos (softwares) utilizem esses recursos. Para isso ele é encarregado em fazer comunicação entre componentes de hardware e software, oferecendo uma camada de abstração de nível mais baixo para os recursos que os aplicativos devem controlar para executar suas ações. 

\subsection{Linux}

O linux é um sistema operacional de código fonte livre, sendo assim disponível para qualquer um utilizar e modificar, respeitando os termos de contrato. O linux foi criado baseando-se no sistema operacional Unix, pelo programador finlandês Linus Torvalds. Os sistemas operacionais linux são caracterizados por possuírem o núcleo kernel linux em suas maquinas. 

\subsection{Driver}

 É um desafio fornecer uma definição única e precisa para o termo driver. No sentido mais simples, um driver é um software que permite que o sistema operacional e um dispositivo se comuniquem um com o outro. Portanto, um driver é um software que traduz o que diz um hardware ou um dispositivo para que o computador possa entender. Sem um software de driver, o hardware conectado (por exemplo, uma placa de vídeo ou impressora) não funcionará corretamente.

\subsection{Módulos}

Em computação, um módulo carregável do núcleo, é um arquivo objeto que contém código para estender o núcleo em execução, ou o chamado núcleo base, de um sistema operacional. Os módulos são normalmente usados para adicionar suporte para novos hardwares (como controladores de dispositivos) e/ou sistemas de arquivos, ou para adicionar chamadas de sistema. Quando a funcionalidade fornecida por um módulo não for mais necessária, ela pode ser descarregada com o objetivo de liberar memória e outros recursos.

\subsection{Electronic codebook (ECB)}

O modo mais simples de criptografia é o electronic codebook (ECB). A mensagem é dividida em blocos e cada bloco é criptografado separadamente. A desvantagem deste método é que blocos idênticos de texto plano são criptografados em blocos de texto cifrado idênticos; assim, ele também não oculta padrões de dados. No geral, não oferece uma perfeita confidencialidade de mensagem, e não é recomendado para uso em protocolos criptográficos em geral.
Eis aqui um bom exemplo do nível no qual o ECB pode transformar os padrões de dados de texto simples em texto cifrado.

\subsection{Advanced Encryption Standard (AES)}

AES é uma criptografia de blocos com chave simétrica (cifra de bloco) no caso o AES-128 foi utilizado, dessa forma trabalha com o sistema de blocos de 16 Bytes. É possível utilizar valores de entrada menores sem problemas, mas maiores será necessário dividir em blocos de 16B.

\subsection{Hash}

Funções hash criptográficas são operações matemáticas tendo um algoritmo de uma via, ou seja, é irreversível. É muito usado com senhas da seguinte forma: Primeiramente é gerado um Hash da senha e este, será apenas comparado ao Hash armazenado no destino. Caso os Hash’s sejam iguais, logo a senha é igual. Por exemplo: MD5, SHA-1, SHA-2.

\subsection{SHA-256}

SHA-256 e SHA-512 são funções Hash inovadoras computadas com palavras de 32 bytes e 64 bytes, respectivamente. Eles usam quantidades de deslocamento e constantes aditivas diferentes, mas as suas estruturas são praticamente idênticas, diferindo apenas no número de rodadas.


\section{O que foi desenvolvido}

\subsection{Progama de usuário em C}

O programa em C possue um menu em que o usuário tem a opção de, encriptar uma string, encriptar escrevendo os bytes em hexadecimal, e decriptando utilizando tambem os valores dos bytes em hexadecimal, o programa então faz as escrita no módulo /dev/cryptomodule, inserindo no primeiro byte da string uma letra equivalente a operação a ser realizada, para então poder ler do mesmo módulo. O programa, também, se selecionado para encriptar um string comum, faz a conversão de uma string para um string com os valores originais representados em hexadecimal.

\subsection{Módulo}

O módulo, acionado após uma escrita em /dev/cryptomodule, recebe a string que o usuário proporcionou, e dependendo da primeira letra da string, faz a operação requisitada pelo o usuário, que pode ser, encriptação, desencriptação ou hash da string, e o modulo faz essas operações utilizando a biblioteca crypto.h.


\section{Desafios Encontrados}

Durante o processo de desenvolvimento foram encontrados varios desafios:

\begin{enumerate}
\item Comunicação entre programa de usuário e módulo
\item Trasformação de uma string com caracteres hexadecimais para sua sua representação original
\item Utilização da biblioteca crypto
\item Scatterlist e como utiliza-las
\end{enumerate}

\subsection{Comunicação entre programa de usuário e módulo}

Essa dificuldade foi facilmente superada após estudo de códigos prontos disponibilizados pelo professor.
Utilizamos a função write e read em C, aprendemos sobre essas funções durante a execução do experimento, pois até o momento todos integrantes só havim utilizados as funções fwrite e fread, o conceito é parecido, entretanto fread e fwrite só são utilizados para arquivos.

\vfill

\subsection{Trasformação de uma string com caracteres hexadecimais para sua sua representação original}

Para esse problema utilizamos manipulação de bits com a operação shifleft \textbf{\textless\textless}, e com com a subtração do caracter para o valor que ele representa utilizando os valores da \textit{ASCII Table}.
\begin{flushleft}
\begin{em}
Ex.:\\
4A      $\Rightarrow$ 00110100 01000001 : 2 caractéres\\
4Ah     $\Rightarrow$ 00000100 00001010 : Caractéres subtraídos para o valor original em hexadecimal\\
400Ah   $\Rightarrow$ 01000000 00001010 : O primero byte é shiftado 4 vezes\\
J / 4Ah $\Rightarrow$ 00101010          : Os dois bytes são somados, gerando o valor original que os dois caractéres representavam \\ 
\end{em}
\end{flushleft}

\subsection{Utilização da biblioteca crypto}
Esse problema não foi específico a uma função ou a uma lógica, mas sim a diferentes mecânismos necessários para a utilização do modulo que o tornou complicado de utilizar, foi possível sua utilização atravez de estudos dos códigos prontos e vários testes. 

\subsection{Scatterlist e como utiliza-las}
Scatterlist é uma struct em que são armazenadas \textit{pagenumber}, \textit{offset} e \textit{size} que são todas as informções para encontrar uma variável na memória, esse método de armazenar os dados de uma variável é muito utilizada em módulos de kernel.\\
Para resolver o problema de encontrar essa variável foi encontrado a função \textit{sg\_virt} que calcula o endereço da variável dando sua respectiva scatterlist, assim possibilitando a utilização dessa variável
\vfill
\clearpage

\section{Conclusão}
Resultados dos experimentos
\subsection{Programa de Usuário em C}
No programa de usuário foi criado um menu com cinco opções:
\begin{enumerate}
    \item Criptografar uma string.
    \item Criptografar uma string em hexadecimal
    \item Decriptografar uma string em hexadecimal
    \item Gerar hash de uma string em hexadecimal
    \item Sair
\end{enumerate}
\textbf{Criptografar uma string:}
Essa opção faz com que o programa converta a string em ASCII para sua conversão em hexadecimal.
E, então, escreve a string convertida, com o primeiro caracter 'c', no módulo \emph{/dev/cryptomodule}.
\\
\\
\textbf{Criptografar uma string em hexadecimal:}
Aqui o programa apenas escreve a string dada pelo usuário o módulo, também com o primeiro caracter 'c'.
\\
\\
\textbf{Decriptografar uma string em hexadecimal:}
Nessa opção também é feita a escrita, entretanto com 'd' como primeiro caracter.
\\
\\
\textbf{Gerar hash de uma string em hexadecimal:}
Aqui a escrita é feita direta, com um 'h' na primeira posição da string.
\\
\\
\textbf{Sair:}
Sair do programa sem realizar nenhuma operação.
\\
\\
Em todas as opçõe, menos \emph{Sair}, o resultado é impresso no final da execução.

\subsection{Módulo cryptomodule}
No módulo, utilizando várias funções específicas de Kernel, a primeira coisa que o módulo faz é reverter a string de entrada, escrita pelo programa de usuário, então a string é tratada adicionando \emph{Padding} quando necessário.
Depois do tratamento, utilizando o primeiro caracter para designar a função, direcionamos a string tratada à função certa equivalente a opção escolhida pelo usuário.
\\\\
O primerio problema que não foi possível contornar é quando a string dada pelo o usuário tem um tamanho real maior que 16 bytes. Já que esse é o tamanho do nosso bloco de encriptação, portanto não foi resolvido a necessidade de multiplas conversões.
\\
\\
Após à execução das funções necessárias o resultado delas é escrito novamente no módulo, para que o programa de usuário possa ler os resultado das operações.


\end{document}
